\section{Parallel Patterns}

\subsection{Divide and Conquer}

\begin{definition}
\textbf{Divide and Conquer } refers to the strategy of splitting a problem into smaller subproblems,
                solving them recursively, and combining the results.
\end{definition}

Problems of DC with manual fork/join:
\begin{itemize}
    \item Poor reusability of code
    \item Commit to a number of threads a priori and can not dynamically adjust.
    \item Subproblems are not automatically assigned to a thread.
\end{itemize}

\begin{remark}
Some manual fork/join fixes are: sequential-cutoffs or running one
subproblem on the thread itself, instad of creating two new threads.
\end{remark}

\subsubsection{Executor Service}

The Executor Service framework provides the way to schedule and run tasks on a thread-pool automatically.

\begin{definition}
\textbf{Runnable}: Task without return value. (Passed to \texttt{execute()} or \texttt{submit()}).
\textbf{Callable}: Task returning \texttt{Future<T>}. (Passed to \texttt{submit()}).
\end{definition}

\begin{remark}
\textbf{Pool Types}: \texttt{newFixedThreadPool(n)} (keeps $n$ threads), \texttt{newCachedThreadPool()} (scales dynamically).
\end{remark}

Pros and Cons of Executor service:
\begin{itemize}
    \item Not suited for: Recursive subproblems that depend on each other and have a deep fork-tree. (run out of threads)
    \item Suited for: flat structures that can run independently in parallel.
\end{itemize}

\subsubsection{Fork-Join}

Manages a pool of threads to solve recursive tasks using a work-stealing strategy, where threads have local task-queues
which they work on and if empty steal from other threads.

\begin{definition}
\textbf{RecursiveTask}: Returns a result (like Callable).
\textbf{RecursiveAction}: Does not return a result (like Runnable).
\textbf{Key Methods}: \texttt{fork()} (async push), \texttt{join()} (wait, but steals work while waiting), \texttt{compute()} (main logic).
\end{definition}

\begin{codeexample}
class MyTask extends RecursiveTask<Integer> {
    int start, end;
    MyTask(int start, int end) { this.start = start; this.end = end; }
    
    protected Integer compute() {
        if (end - start <= SEQUENTIAL_CUTOFF) { 
            return baseCase(); 
        } else {
            int mid = (start + end) / 2;
            MyTask left = new MyTask(start, mid);
            MyTask right = new MyTask(mid, end);
            
            left.fork(); // Async execute left half
            int rightRes = right.compute(); // Optimize: compute right half now
            int leftRes = left.join();      // Wait for left half
            
            return leftRes + rightRes;
        } 
    }
}
\end{codeexample}

\begin{important}
Make sure that every task has enough work to do (aprox. 100-10'000 ops.) (else the overhead will be larger than the parallel speedup gain)
\end{important}

\subsubsection{Task graphs}

The execution of parallel programs using Fork-Join can be represented as a Task-Graph (DAG).

\begin{definition}
$T_1$ is the execution time on a sequential machine (sum of all task durations in the DAG).
\textbf{Span} $T_{\infty}$: execution time with unlimited hardware (sum of longest dependent path in DAG).
Max speedup is limited by $T_1 / T_{\infty}$ (Amdahl's Law).
\end{definition}

\begin{remark}
It is the programmers responsibility to find the optimal decomposition for minimal span (balanced tree). \textbf{Fork-join guarantees optimal scheduling with minimal idle time.}
\end{remark}

\subsection{Parallel Patterns}

\begin{definition}
A \textbf{reduction} produces a single answer from a collection via an associative operator.
\end{definition}

\begin{example}
Max, Min, Sum, Product of an array.
\end{example}

\begin{definition}
A \textbf{map} applies a function to each elements of a collection producing an output of the same size.
\end{definition}

\begin{example}
Vector addition: $v_i = u_i + w_i \forall i \in [1, n]$
\end{example}

\begin{remark}
Maps and Reductions have an optimal span of $O(\log n)$.
\end{remark}

\begin{definition}
A \textbf{stencil} extends the map pattern by using the element itself and its neighbours.  
\end{definition}

\begin{example}
A smoothing kernel (ex. 3x3) that calculates the average value is sled over an image to smooth/blur the image.
\end{example}

\begin{remark}
Arrays and balanced trees are usually parallelizable while linked lists are not even if the per element operation is parallelizable.
\end{remark}

\subsection{Parallel Algorithms}

\subsubsection{Prefix-Sum}

Prefix-Sum: Sum of first $n$ inputs.

\begin{algorithm}
upTask(); //divide problem and build tree with sums of subproblems
downTask(); //fill in the correction values (left_sum + parent_correction)
\end{algorithm}

\begin{theorem}
The span of the prefix sum optimization is $O(\log n)$, therefore parallelism of $\frac{n}{\log n}$.
\end{theorem}

\subsubsection{Pack}

\textbf{Pack} pack all elements from an array that satisfy a certain condition into a seperate output. (by using prefix sum)

\begin{algorithm}
applyMap(); //applies a map to the input
//results in a binary vector [1, 0, 1, 1, 0, ...]\\
parallelPrefixSum(); //calculates the prefix sum of the binary vector\\
//results in [1, 1, 2, 3, 3, ...]\\
applyPackMap(); //applies another map to the result of the prefix sum\\
//results in [0, 2, 3, ...] (indices of elements to write to output)
\end{algorithm}

\subsubsection{Quicksort}

We can reduce the \textbf{limiting factor of the problem division} from a span of $O(n)$ to a span of $O(\log n)$
using the pack-pattern.

\begin{algorithm}
applyMaps(); //apply maps for greater and smaller than pivot
parallelPrefixSum(); //count how many smallerthan/greatherthan elements come bofore
applyMap(); //index = numLessThan || numTotalLessThan + numGreaterThan
\end{algorithm}

We can do this by using two pack-patterns to partition the data. This leads to a partitioning of $O(\log n)$
levels times $O(\log n)$ for prefixSum hence $O(\log^2 n)$ total span. 

This means a parallelism of $\frac{n\log n}{\log^2 n} = \frac{n}{\log n}$.
